\section{Proposition de budget}

 > Le tableau suivant présente une description point par point du budget demandé ainsi que de sa destination. Les principaux postes budgétaires correspondent à la rémunération des mentors, aux rafraîchissements pour les rencontres du programme et, enfin, à un budget pour une expérience collective visant à renforcer la cohésion du groupe. Les coordonnateurs ont établi ces estimations à partir de reçus et de tarifs réels. Nous avons également prévu une marge pour imprévus.

\medskip

\renewcommand{\arraystretch}{1.4}
\noindent\begin{tabularx}{\textwidth}{|X|>{\centering\arraybackslash}p{0.2\textwidth}|>{\centering\arraybackslash}p{0.25\textwidth}|}
\hline
\rowcolor{uqamblue!85}
\textcolor{white}{\textbf{Poste budgétaire}} & \textcolor{white}{\textbf{Montant estimé}} \\
\hline
Rémunération des mentors & 2100 CAD \\
\hline
Rencontres mensuelles (collations) & 500 CAD\\
\hline
Activité de groupe (à définir) & 250 CAD\\
\hline
Fonds de prévoyance & 150 CAD\\
\hline
\textbf{Total} & 3000 CAD\\
\hline
\end{tabularx}
\renewcommand{\arraystretch}{1}